\section{The object: metallic once-punctured-torus bundles}
\label{sec:object}

Let $T_\circ$ be the once-punctured torus and $\mathrm{MCG}(T_\circ)\cong\SL(2,\Z)$ its
mapping class group, generated by the two Dehn twists $R$ and $L$ along a dual pair of
simple closed curves. To a pseudo-Anosov word $\varphi$ in $R$ and $L$ one associates
the mapping torus $M_\varphi = T_\circ\times[0,1]/(x,1)\sim(\varphi x,0)$, a cusped
hyperbolic three-manifold fibring over the circle. Our object is the subfamily whose
monodromies are the \emph{metallic words}
\begin{equation}
  W_m \;=\; R^m L^m,\qquad m=1,2,3,\dots
\end{equation}
together with their cusp-glued composites $R^{m_1}L^{m_1}\cdots R^{m_k}L^{m_k}$, which we
return to in Section~\ref{sec:chirality}.

\subsection{Why ``metallic''}

In the standard representation $R\mapsto\left(\begin{smallmatrix}1&1\\0&1\end{smallmatrix}\right)$,
$L\mapsto\left(\begin{smallmatrix}1&0\\1&1\end{smallmatrix}\right)$ of
$\mathrm{MCG}(T_\circ)$, a single metallic block is conjugate to a power of the matrix
\begin{equation}
  M_m \;=\; \begin{pmatrix} m & 1 \\ 1 & 0 \end{pmatrix},
  \qquad \det M_m = -1,\qquad \tr M_m = m,
\end{equation}
whose larger eigenvalue is the \emph{metallic mean}
$\lambda_m=\tfrac12\!\left(m+\sqrt{m^2+4}\right)$, the positive solution of the
self-referential equation $x = m + 1/x$. For $m=1$ this is the golden mean
$\varphi=\tfrac12(1+\sqrt5)$; for $m=2$ the silver mean $1+\sqrt2$. The dilatation of the
pseudo-Anosov monodromy $W_m$ is $\lambda_m^2$; the larger $m$, the more strongly the
fibre is stretched (\fig{fig:metallic}).

\begin{figure}[t]
  \centering
  \includegraphics[width=\linewidth]{fig01_metallic_means}
  \caption{The metallic seeds. \emph{Left:} the metallic mean $\lambda_m$ and the
  monodromy dilatation $\lambda_m^2$ as functions of the block parameter $m$; the
  golden ($m=1$, figure-eight) and silver ($m=2$) cases are marked. \emph{Right:} the
  self-referential recursion $x=m+1/x$ converging to $\varphi$ --- the object literally
  generated by iterating a minimal self-reference. \sProved}
  \label{fig:metallic}
\end{figure}

The determinant $\det M_m=-1$ is not incidental; it is the source of the
orientation-reversing structure that will govern chirality in Section~\ref{sec:chirality}
and the alternating sign $(-1)^{n-1}$ in the $A$-polynomial family of
Section~\ref{sec:knot}. We record it as the object's first structural fact.

\begin{observation}[\sProved]
The eigenvalues of $M_m^{-1}$ coincide with those of $-M_m$. Indeed if $\lambda,\mu$ are
the eigenvalues of $M_m$ then $\lambda\mu=\det M_m=-1$, so $\lambda^{-1}=-\mu$ and
$\mu^{-1}=-\lambda$; hence $\chi(M_m^{-1})=\chi(-M_m)$ as characteristic polynomials.
This identity is why the ``inverse'' and ``sign'' directions of the tower
(Section~\ref{sec:tower}) coincide.
\end{observation}

\subsection{The figure-eight as the golden case}

The case $m=1$ is the figure-eight knot complement $4_1$. Its monodromy is
$\varphi=R L$, conjugate to $M_1^2=\left(\begin{smallmatrix}2&1\\1&1\end{smallmatrix}\right)$,
and it is the unique arithmetic once-punctured-torus bundle of smallest volume,
$\Vol(4_1)=2.029883\ldots=2V_{\mathrm{tet}}$, with invariant trace field $\Q(\sqrt{-3})$.
It is amphichiral (it admits an orientation-reversing self-homeomorphism), and its
Chern--Simons invariant vanishes. These two facts --- minimality and amphichirality ---
make the figure-eight the natural \emph{base} of the family, and, as we will see in
Section~\ref{sec:bridge}, the object \emph{least} able to carry a physical scale. The
silver case $m=2$ is the bundle with monodromy $R^2L^2$, trace field $\Q(i)$, the second
and last arithmetic member of the canonical metallic series.

\subsection{Trace fields and the arithmetic backbone}

The monodromy matrix $M_m$ has real eigenvalues lying in the quadratic field
$\Q(\sqrt{m^2+4})$. This is the field of the \emph{dynamics} --- the dilatation field ---
and is distinct from the invariant trace field of the hyperbolic structure on $M_{W_m}$,
which is in general a complex field (e.g.\ $\Q(\sqrt{-3})$ for the figure-eight). Keeping
these two fields apart is a recurring discipline of the subject; we use the dilatation
field $\Q(\sqrt{m^2+4})$ whenever we speak of the trace-map dynamics, and name the
invariant trace field explicitly when we mean the hyperbolic geometry. The arithmetic
backbone of the canonical family is summarised in Table~\ref{tab:fields}.

\begin{table}[h]
\centering\small
\begin{tabular}{cllll}
\toprule
$m$ & word & metallic mean & dilatation field & note \\
\midrule
$1$ & $RL$ & golden $\varphi$ & $\Q(\sqrt5)$ & figure-eight $4_1$; trace field $\Q(\sqrt{-3})$ \\
$2$ & $R^2L^2$ & silver $1+\sqrt2$ & $\Q(\sqrt8)$ & ``silver'' bundle; trace field $\Q(i)$ \\
$m$ & $R^mL^m$ & $\tfrac12(m+\sqrt{m^2+4})$ & $\Q(\sqrt{m^2+4})$ & generic metallic bundle \\
\bottomrule
\end{tabular}
\caption{The arithmetic backbone of the canonical metallic family. \sProved}
\label{tab:fields}
\end{table}
